% 本文件由 LaTeX 排版 Agent 根据有效 translation.json 自动生成。
% author=Tertullian; work_id=0401; title=论披肩

\chapter{论披肩}

% structure: p0001 source=h1 target=omitted reason=page-title-already-rendered-by-parent

% structure: p0002 source=h2 target=section reason=relative-heading-rank; base=h2

\section{第一章  时移世易，服饰与命运俱变}

\Place{迦太基}人啊，永远是\Place{阿非利加}的君王，因古昔的记忆而高贵，因今日的福祉而蒙福。我欣喜你们时代如此繁荣，竟有闲暇和乐趣来批评服饰。这正是\Quote{太平盛世}和丰裕的时期。福泽从天朝和天穹降下。然而，你们古时也穿着另一形状的服装——你们的常服。它们确实因织艺精巧、色调和谐、尺寸合宜而受称誉：既不过分长至胫骨，也不失礼地短于膝间；既非吝啬于臂，亦非紧束于手；而且，没有腰带分割衣褶的遮掩，它们以方正对称之态披在男子背上。外套上的披风——本身亦为方形——搭在双肩，在颈处由搭扣紧合，曾安放于肩头。其对应物如今是司祭的礼服，敬献给\Person{阿斯克勒庇俄斯}，你们如今宣称他属于你们。同样，在你们近邻，姊妹城邦曾以此装束其公民；还有在\Place{阿非利加}其他推罗人定居之地。但当日世命运的签简改变，\Person{天主}眷顾\Place{罗马}人时，那姊妹城邦确实自愿仓促变革；当\Person{西庇阿}停靠其港口时，她已提前以服饰向他致意，在罗马化上早熟。至于你们，在你们遭受损害而产生的利益之后，既豁免了时间的无限，也未剥夺你们崇高的地位——在\Person{格拉古}及其不祥征兆之后，在\Person{李必达}及其粗鲁玩笑之后，在\Person{庞培}及其三重祭台之后，在\Person{凯撒}及其长期拖延之后，当\Person{斯塔提利乌斯·托鲁斯}筑起你们的城墙，\Person{森提乌斯·萨图尔尼努斯}宣读你们建城的庄严仪式时——正当和睦伸出援手，那\Emph{长袍}被献上。啊！它走了多么远的历程！从佩拉斯吉人到吕底亚人；从吕底亚人到罗马人：为的是从更崇高民族的肩头降下，拥抱\Place{迦太基}人！从此，你们觉得常服太长，便用腰带将其束起；你们将如今平滑的\Emph{长袍}的多余部分叠聚起来支撑；无论社会地位、尊严或季节让你们穿上何种其他衣物，那\Emph{披风}——曾被你们各阶层、各处境者所穿——你们不仅遗忘，甚至嘲笑。至于我，面对一个更古老的证据（你们的遗忘），我并不惊讶。因为那撞锤——并非\Person{拉贝里乌斯}所言的——\par

\Quote{角向后扭、毛皮为肤、曳石而行}\par

但它是一种梁状机具，用于撞墙的军役——从未被任何人使用过的，令人敬畏的\Place{迦太基}，\par

\Quote{最锐于征战之事}\par

据说（\Place{迦太基}）是第一个装备了用于摆动冲力的振荡工作的；将其引擎的力量比照那报仇之兽的狂怒而塑造。然而，当他们的国家命运垂危，那撞锤——如今已变成\Place{罗马}的——正对着曾属于它自己的城墙大显神威时，迦太基人立刻目瞪口呆，如同见到一种\Quote{新奇}和\Quote{陌生}的巧技：\Quote{漫长岁月竟有如此变迁之力！}简言之，披风也是如此，不再被识别了。\par

% structure: p0008 source=h2 target=section reason=relative-heading-rank; base=h2

\section{第二章  变化的法则，或普世变异}

现在让我们从别处取材，免得布匿特性在罗马人中要么脸红，要么悲伤。无论如何，改变习惯是整个自然界既定的功能。就连我们所居住的这个世界本身，也在此期间执行着这一点。阿那克西曼德，如果他以为还有更多的世界，就让他去留意吧：还有任何其他人，如果以为在墨洛普斯区域的某处存在着另一个世界，正如西勒诺斯在迈达斯耳边胡言乱语（那双耳朵确实适合听更大的神话）那样，就让他去留意吧。不仅如此，即使柏拉图以为存在一个世界，而我们的世界是它的影像，那个世界也必然同样要经历变化；因为，如果它是一个\Quote{世界}，它就会由不同的实体和功能构成，与在这里被称为\Quote{世界}的形式相称：因为如果它不正好与\Quote{世界}相同，它就不是\Quote{世界}。那些在多样性中趋向统一的事物，是\Emph{通过变异}而多样的。简言之，正是它们的变迁使它们多样性的不和谐得以联合。因此，每一个\Quote{世界}都将\Emph{通过变化}而存在，其整体结构源于多样性，其调和源于变迁。无论如何，我们这客栈是\OriginalTerm{变化多端}{versiform}的——这一点对于闭着的眼睛，或者完全\Person{荷马}式的眼睛，也是显而易见的。白昼和黑夜轮流交替。太阳随季节变化，月亮随月相变化。群星——在纷乱中各具特色——有时隐没，有时又几乎复明。天穹时而晴朗灿烂，时而阴云密布；或者暴雨倾泻，连同各种混合的投射物；此后是微微细雨，然后又放晴。同样，大海在诚实方面名声不佳；有时，微风均匀地吹拂，平静赋予它正直的表象，风平浪静赋予它温和的表象；然后突然之间，它波涛汹涌，如山峦起伏。同样，如果你观察大地，它喜欢适时地装扮自己，你几乎会否认它的同一性，当你想起它的绿色，却又看见它金黄，很快还会看见它灰白。至于它其余的装饰，有什么不是处于交替变化之中的呢？——山脉较高的山脊因冲刷而降低，泉源的水脉因消失而干涸，河流的通道因冲积而改变。曾有一个时期，它的整个球体都被水淹没，经历了变化。至今，海贝和三叉螺号仍作为异乡客留居在山上，渴望向\Person{柏拉图}证明，就连高地也曾波动过。但是，随着退潮，它的球体又经历了一次形式上的变化；另一个，却是同一个。甚至现在，它的形状还在经历局部变化，当某个地点受损时；当在它的岛屿中，\Place{德洛斯岛}如今已不复存在，\Place{萨摩斯岛}成了一堆沙，\Person{西比尔}并非说谎者；当在\Place{大西洋}中，那片曾与\Place{利比亚}或\Place{亚洲}同样大小的岛屿已无处可寻；当从前\Place{意大利}的一侧，被亚得里亚海和第勒尼安海的震动断裂到中心，留下\Place{西西里}作为其遗迹；当那次完全的分裂，使海洋的冲突向后旋转，赋予大海一种新的恶习，不是吐出残骸，而是吞噬它们！大陆也遭受来自天上或内在的力量。看看\Place{巴勒斯坦}。那里，\Place{约旦}河是边界的裁决者，（看哪）一片广袤的荒芜，一片被剥夺的地区，一片无用的土地！从前那里曾有城市，繁荣的民族，土地出产果实。后来，既然天主是审判者，不敬虔招致了火雨：\Place{索多玛}的日子结束了，\Place{哈摩辣}不复存在；一切都成了灰烬；邻近的大海和土地一样，经历着活死状态！这样的云笼罩了\Place{埃特鲁里亚}，烧毁了它古老的\Place{沃尔西尼}，以此教训\Place{坎帕尼亚}（更因\Place{庞贝}的被夺而）期待自己的山脉。但希望这样的灾难不再重复！但愿\Place{亚洲}此时也无需担心土地的贪婪！但愿\Place{非洲}已经在那吞噬一切的深渊面前颤抖过一次，并以一个营地的奸诈吞噬而赎罪！此外还有许多其他这样的损害，改变了我们球体的形态，并移动了其中的特定地点。战争的放纵也是非常巨大的。但叙述悲惨的细节并不比叙述王国的变迁更令人厌烦，以及（表明）它们的变化有多么频繁，从\Person{尼诺斯}——\Person{贝洛斯}的后代——开始；如果\Person{尼诺斯}确实是第一个拥有王国的人，正如古代异教权威所断言的那样。在他之前，笔通常不会在你们异教徒中旅行。从\Place{亚述人}开始，也许\Quote{有记载的时间}的历史开始展开。然而，我们这些习惯阅读\Emph{神圣}历史的人，从宇宙本身的诞生起就掌握这个主题。但我此刻更喜欢\Emph{愉快}的细节，因为愉快的事物也会变化。简言之，无论海洋冲刷了什么，天空烧毁了什么，大地掏空了什么，利剑砍伐了什么，都会在补偿的轮回中在另一个时间重新出现。因为在最初的日子里，不仅大地在其大部分区域是空旷无人居住的；而且如果某个民族占据了一部分，它也是独自存在的。于是，（大地）终于明白万物都崇拜自己，便打算清除和刮去其（居民的）丰裕，在一个地方密集聚集，在另一个地方放弃它们的驻地；以便从此（如同从接穗和插条）民族从民族，城市从城市，可以在其球体的每一个区域里栽种。过剩民族的蜂群带来了迁徙。\Place{斯基泰人}的富余使\Place{波斯人}肥沃；\Place{腓尼基人}涌入\Place{非洲}；\Place{弗里吉亚人}生出\Place{罗马人}；\Place{迦勒底人}的种子被带进\Place{埃及}；随后，当从那里转移时，它成为\OriginalTerm{犹太的}{Jewish}种族。同样，\Person{赫拉克勒斯}的后裔同样为了\Person{泰梅诺斯}的利益前去占据\Place{伯罗奔尼撒}。同样，\Person{涅琉斯}的\Place{爱奥尼亚}伙伴为\Place{亚洲}造就了新城市；同样，\Person{阿尔基亚斯}领导的\Place{科林斯人}加强了\Place{叙拉古}。但如今提及古物已是徒劳，因为我们自己的进程就在眼前。这个时代改造了我们球体的多大一部分！我们现存帝国的三重力量产生、或扩大、或恢复了多少城市！当天主眷顾如此众多的\Person{奥古斯都}团结一致时，多少人口被转移到其他地方！多少民族被削减！多少等级被恢复到古代的辉煌！多少蛮族被挫败！的确，我们的球体是这个帝国巧妙耕种的产业；每一个敌意的乌头都被根除；暗中狡猾亲近的仙人掌和荆棘被完全拔起；（球体本身）比\Person{阿尔基努斯}的果园和\Person{迈达斯}的玫瑰园更令人愉悦。因此，既然赞美我们球体的变化，为什么你们却对一个人指指点点呢？\par

% structure: p0010 source=h2 target=section reason=relative-heading-rank; base=h2

\section{第三章 同样受变化法则支配的走兽}

走兽也一样，它们改变的不是\Emph{衣服}，而是\Emph{外形}。然而孔雀仍然拥有羽毛作为衣服，而且是最精美的衣服；不仅如此，它颈部的盛开胜过任何紫色，背部的光辉比任何镶边更金光闪烁，尾巴的拖曳比任何裙裾更飘逸；颜色繁多，多样，变幻；从未是自身，却一直是另一个，尽管当是另一个时又是自身；总之，变化如同移动一样频繁。蛇也值得一提，尽管与孔雀不可同日而语；因为它也完全改变了赋予它的东西——它的皮和年龄：如果确实如此（事实正是这样），当它感到衰老遍及全身时，它将自己挤进狭窄之处；爬进洞穴，同时从皮中爬出；并在当场被剪得干干净净，一越过门槛就立刻将蜕皮留在那里，然后以新的青春展开身体：随着鳞片，它的岁月也被抛弃了。鬣狗，如果你观察，性属一年一变，交替雌雄。我不说牡鹿，因为它自己，作为自己年龄的见证，以蛇为食，因毒的效果而衰弱——进入青春。此外，还有，\par

一种缓行的、田野出没的四足兽，\newline{}谦卑而粗野。\par

你以为是指\OriginalTerm{帕库维乌斯}{Pacuvius}笔下的乌龟吗？不是。还有另一只小兽适合这首短诗；体型非常适中，但名字却很夸张。如果你事先不认识它，听人说\OriginalTerm{变色龙}{chameleon}，你会立刻联想到某种与狮子结合的更巨大的东西。但当你遇到它时（通常在葡萄园里），它的整个身体都藏在一片葡萄叶下，你会立刻嘲笑这个名字的极度狂妄，因为它的体内甚至没有水分，尽管在远更微小的生物中，身体\Emph{是}液化的。变色龙是一张活的薄膜。它的头直接从脊椎开始，因为它没有脖子：因此它很难反思；但在观察上，它的眼睛是突出的，不，它们是旋转的光点。迟钝而疲惫，它几乎不抬离地面，而是拖着茫然的脚步，向前移动——它更多是展示而非迈出一步：永远禁食，却从不昏厥；张着嘴进食；像风箱一样起伏，反刍；它的食物是风。然而，变色龙能够实现完全的自我变化，仅此而已。因为，虽然它的颜色本来是一种，但每当有什么东西接近它时，它就变红。唯独变色龙被赐予——正如我们常说的——用自己的皮玩耍。\par

说了这么多，是为了经过充分准备，我们终于到达\Emph{人}。无论你承认他起源于何处，总之他从造物主的手中出生时是赤裸的、未着衣的：此后，他最终未经许可，就用早熟的抓取占有了智慧。于是急忙遮盖他新造的身体中尚未因羞耻而应遮盖的部分，他暂时用无花果叶围住自己：后来，因犯罪被逐出出生地后，他穿着皮衣，走向世界如同走向矿场。\par

但这些是奥秘，并非所有人都能知晓。来，让我们听听你仓库里的故事——\OriginalTerm{埃及人}{Egyptians}讲述、\OriginalTerm{亚历山大}{Alexander}整理、他母亲阅读的故事——关于\OriginalTerm{奥西里斯}{Osiris}的时代，那时拥有丰富羊群的\OriginalTerm{阿蒙}{Ammon}从\OriginalTerm{利比亚}{Libya}来到他那里。简言之，他们告诉我们，\OriginalTerm{墨丘利}{Mercury}在他们中间时，偶然抚摸了一只公羊的柔软，感到欣喜，便剥了一只小母羊的皮；然后，他坚持不懈地尝试，并由于材料的柔韧，通过持续的拉伸将线变细，织成了他之前用亚麻条编织的原始网形状。但\Emph{你们}宁愿将羊毛加工和织机结构的所有管理都归于\OriginalTerm{弥涅耳瓦}{Minerva}；而\OriginalTerm{阿拉克涅}{Arachne}掌管着一个更勤奋的作坊。从此材料丰富。我不提\OriginalTerm{米利都}{Miletus}、\OriginalTerm{塞尔格}{Selge}和\OriginalTerm{阿尔提努姆}{Altinum}的羊，也不提\OriginalTerm{塔伦图姆}{Tarentum}或\OriginalTerm{贝提卡}{Bætica}著名的羊，它们以自然为染料：但我提到的事实是，灌木丛为你提供衣服，亚麻的草质部分经过洗涤失去绿色而变白。且\Emph{种植}和\Emph{播种}你的外衣还不够，除非你还幸运地\Emph{捕捞}衣服。因为大海也产羊毛，因为更光鲜的苔状羊毛外壳提供了毛状物质。此外，众所周知，\OriginalTerm{蚕}{silkworm}——一种小蠕虫——她安全地重新生产出（绒毛线），通过将线穿过\Emph{空气}，她比蜘蛛的日晷状网更巧妙地拉伸它们，然后吞噬。同样，如果你杀死它，你缠绕的线立即充满鲜艳的颜色。\par

因此，裁剪艺术的巧思，叠加并追随如此丰富的材料库存——首先出于对人性的渴求，必要引领道路；随后也出于装饰，甚至膨胀，野心随之而来——已经颁布了各种形式的服装。这些形式中，部分由特定民族穿着，而不为其他民族所共有；部分则普遍适用，对所有人都有用：例如，这件\OriginalTerm{斗篷}{Mantle}，虽然更希腊化，但如今在言语上已在\OriginalTerm{拉丁姆}{Latium}找到了家。随着词语，服装也传入。因此，那位曾经将\OriginalTerm{希腊人}{Greeks}逐出城市，但到了年老时却学习了他们的\Emph{字母}和\Emph{语言}的人——正是这个\OriginalTerm{加图}{Cato}，在他担任执政官期间裸露肩膀，通过他斗篷般的\Emph{装束}同样地表达了对希腊人的青睐。\par

% structure: p0017 source=h2 target=section reason=relative-heading-rank; base=h2

\section{第四章 变化并非总是改善}

怎么，如今，如果罗马风尚对人人都是（社会的）救恩，你们却为何在某种程度上是希腊人，甚至在那些不光彩的事情上也是如此？或者，如果不这样，那么世界上那些受过更好教养的省份，那些天性更适于通过艰苦努力克服土壤困难的省份，又从哪里得来角力场的追求——那些追求陷入可悲的老迈并徒劳无功——还有泥浆的涂抹、沙中的翻滚、干燥的饮食？从哪里得来，我们有些努米底亚人，他们用马尾羽饰使长发变得更长，学会了让理发师将他们的皮肤刮得干干净净，唯独豁免其头顶不受刀剪？从哪里得来，毛发浓密而粗犷的人学会了教树脂如此贪婪地啃食他们的手臂，拔毛钳如此偷窃般地梳理他们的下巴？真是奇事，这一切竟是在没有曼特行袍的情况下完成的！这整个亚洲式的做法都属于曼特行袍！利比亚啊，还有你，欧洲啊，你们与那些你们不知如何装束的体育上的讲究有何相干？因为，事实上，操练希腊式的脱毛比操练希腊式的服饰算得上什么呢？\par

服饰的转移接近罪责，其程度恰如遭受改变的不是习俗，而是本性。时间应得的尊荣与宗教之间有着足够大的差异。让习俗忠于时间，本性忠于天主。因此，拉里萨的英雄给了本性一击，变成了处女；他是在野兽的骨髓中养大的（由此也得出他名字的组成，因为他用嘴唇与母亲的乳房是陌生的）；他是在一处岩石遍布、林栖和怪异的教练手下，在一所石头学校里长大的。你会耐心忍受，如果是在一个\Emph{男孩}的例子中，那是他母亲的操心；但无论如何，他已经长了头发，无论如何他已经秘密地向某人证明了他的成年，当他同意穿上飘逸的袍子、梳理头发、保养皮肤、照镜子、装饰脖颈时；甚至因穿孔而变得女性化，连耳朵也在其中，\Place{西革翁}的半身像仍保留着那痕迹。后来他显然成了士兵：因为需要恢复了他的性别。战斗的号角已经吹响：武器也并不难寻。\Quote{钢铁本身，}（\Person{荷马}）说，\Quote{吸引英雄。}否则，如果在那激励之后如同之前一样，他坚持他的处女身份，他或许也可能结婚了！看，因此，突变！我称他为怪物——双重的怪物：从男人变为女人；不久又从女人变回男人：然而真理本不应被歪曲，欺骗也不应被承认。每种变化的方式都是邪恶的：一种违背本性，另一种违背安全。\par

然而，当情欲在衣着上改变一个人时，其羞耻更甚于某种母亲的恐惧所导致的改变：但你们却向我献上敬拜，你们本应为我羞愧——那位持棍持毛皮者，他用女性的装束换掉了自己整个骄傲的名字遗产！这种放纵被允许给\Place{吕底亚}的隐秘角落，以致\Person{赫拉克勒斯}在\Person{翁法勒}身上被卖淫，\Person{翁法勒}也在\Person{赫拉克勒斯}身上被卖淫。\Person{狄俄墨得斯}和他血腥的马槽在哪里？\Person{布西里斯}和他丧葬的祭台在哪里？\Person{革律翁}，三重一体在哪里？棍棒仍然更喜欢以他们的脑浆为渍，当它被香膏烦扰时！那陈年（污渍）的\Person{许德拉}血和半人马血在箭杆上渐渐被浮石磨掉，浮石熟悉发簪！而享乐侮辱了这一事实：在刺穿怪物之后，它们或许会缝制一个冠冕！没有一位正派的女人，甚至任何著名的女英雄，会冒险将她的肩膀放在这样一头野兽的毛皮之下，除非经过长时间的软化、磨平和除臭（我希望在\Person{翁法勒}家里，这是用香脂和胡芦巴膏完成的：我想鬃毛也服从了梳子），以免她娇嫩的脖子沾染上狮子的刚硬。张开的嘴里塞满毛发，颚齿在前额乱发中显得阴暗，整个受辱的面孔，若是能吼叫，早就吼叫了。\Place{涅墨亚}，至少（如果那地方有守护神的话），呻吟了：因为那时她环顾四周，看到她已经失去了她的狮子。这位\Person{赫拉克勒斯}在\Person{翁法勒}的丝绸中变成了什么样子，对\Person{翁法勒}在\Person{赫拉克勒斯}毛皮中的描述已间接描绘了出来。\par

但是，那曾与提任斯人竞争的人——拳击手\Person{克利奥马科斯}——随后在\Place{奥林匹亚}，由于一种令人难以置信的突变，通过流出而失去了他的男性性别——他的皮肤内外青肿，值得在\Person{诺维乌斯}的\WorkTitle{洗衣匠}中戴冠，并理所当然地被哑剧作家\Person{伦图卢斯}在他的\WorkTitle{卡提那人}中纪念——他自然不仅用臂钏遮盖了拳击带留下的痕迹，而且用某种最精细的织物替代了他运动员斗篷的粗糙和粗犷。\par

关于\Person{菲斯可}和\Person{萨丹纳帕路斯}，我必须保持沉默，若不是他们在情欲中的卓越，没人会认出他们是国王。但我必须保持沉默，以免甚至\Emph{他们}也会低声议论你们某些恺撒，同样不知羞耻；以免命令已给予犬类的忠诚，去指着一个比\Person{菲斯可}更不洁、比\Person{萨丹纳帕路斯}更软弱，而且确实是第二个\Person{尼禄}的恺撒。\par

虚浮光荣的力量也同样热切地促使服装的变迁，即使男子气概得以保存。每一种情感都是一种热；然而，当它被吹到（火焰般的）矫揉造作时，立刻因光荣的烈焰而成为热情。因此，从这燃料中，你看见一位伟大的君王——仅次于他的光荣——在沸腾。他征服了米底亚民族，却被米底亚服装所征服。脱去胜利的铠甲，他降格为俘虏的裤子！那饰有鳞状浮雕的胸膛，通过覆盖一层透明织物而裸露；仍然追求战争的作品，（并仿佛）软化，用透气的丝绸将其熄灭！马其顿人的精神还不够膨胀，除非他也同样喜欢膨胀的服装：只是哲学家们（我相信）自己也有些喜欢那种东西；因为我听说曾经有穿紫色衣服进行哲学思考的。如果一位哲学家穿紫色，为什么不穿镀金拖鞋呢？提洛人除了金鞋之外什么鞋都不穿，这完全不符合希腊习惯。某人会说：\Quote{好吧，但另一个人确实穿丝绸，并穿着\Emph{铜}拖鞋。}确实，值得如此，以便在他那酒神狂欢的服装底部发出一些叮当声，他走在铙钹上！但如果那时\Person{第欧根尼}从他的木桶里吠叫，他不会用泥脚踩他——正如柏拉图式的躺椅所证——而会把\Person{恩培多克勒}直接带到\Place{克娄喀娜}隐秘的深处；以便那个疯狂地以为自己是一个天上存在的人，作为一个神，首先问候他的姐妹们，然后问候男人。因此，这样的服装，由于疏离自然和端庄，我们应当正视、用手指指、点头嘲弄。同样，如果一个人穿着像\Person{米南德}那样娇气的拖地精致长袍，他会听到喜剧演员对他说的：\Quote{那个疯子浪费的是什么样的斗篷？}因为，审查警惕的皱眉头早已抚平，就指责而言，混杂的用法向我们展示了身着骑士服装的被释奴，身着绅士服装的烙刑奴隶，身着自由民服装的声名狼藉者，身着市民服装的乡巴佬，身着律师服装的小丑，身着军装的乡下人；抬棺者、皮条客、角斗士教练，他们都和你一样穿着。再转向女人。你要看到\Person{凯基纳·塞维鲁}严肃提请元老院注意的事——贵妇们在公共场合不穿斯托拉。事实上，占卜官\Person{伦图卢斯}的法令对任何这样放弃斯托拉的贵妇施加的惩罚与通奸相同；因为某些贵妇积极促进废弃那些作为尊严证据和守护者的服装，因为它们是实行卖淫的障碍。但现在，在她们自我卖淫中，为了更容易被接近，她们抛弃了斯托拉、衬衣、帽子、便帽；是的，甚至抛弃了即使在公共场合也用来保持隐私和隐秘的轿子和肩舆。但是，当一个人熄灭了自己合宜的装饰，另一个人却闪耀着不属于自己的装饰。看看街头妓女，大众欲望的屠宰场；还有那些与同性自虐的女人；如果最好从这些公开屠杀贞洁的可耻景象中移开目光，但斜眼一看，你会立刻看到贵妇们！而妓院看守人展示她鼓起的丝绸，用项链安慰她的脖子——比她出没的地方更不洁——并在臂环中插入（即使是贵妇们也毫不犹豫地将勇敢者所得的奖赏据为己有）那些对一切羞耻之事都知情的手，她给她不洁的腿穿上纯白或粉色的鞋子；你为什么不盯着这样的服装呢？或者，再看那些虚假地以宗教作为其新奇支持者的服装？为了白色衣服、头带标志和头盔特权，一些人被引入\Person{刻瑞斯}的秘仪；而由于对深色衣服的相反渴望，以及头上阴郁的羊毛覆盖物，另一些人在\Person{贝娄娜}神庙中疯狂；而用更宽的紫色条纹长袍包裹自己、肩上披加拉太白红色斗篷的吸引力，使\Person{萨图尔努斯}受到爱戴。当这斗篷本身以更严格的用心布置，以及希腊式凉鞋，用来讨好\Person{埃斯库拉庇乌斯}时，你更应当用它来指责和攻击它，认为它犯了迷信的罪——尽管是简单而不造作的迷信？当然，当它第一次穿上这种智慧，即弃绝所有虚荣的迷信时，那么最确定地，斗篷就是一件庄严的袍子，胜过你用来装扮你的男女神祇的所有服装；并且，胜过你的\Person{萨利}祭司和\Person{弗拉米内}司祭的所有帽子和饰带。垂下你的眼睛，我建议你，敬畏这服装，仅凭这一个理由（同时，不等待其他），即它是你错误的弃绝者。\par

% structure: p0024 source=h2 target=section reason=relative-heading-rank; base=h2

\section{第五章 斗篷的美德。它为自己辩护。}

\Quote{然而，}你说，\Quote{我们必须这样从长袍改为斗篷吗？}怎么，如果从王冠和权杖改换呢？当\Person{阿那卡西斯}更喜爱哲学胜过\Place{西徐亚}的王权时，他难道不也是改换了吗？即使没有（奇迹般的）迹象证明你的转变变得更好，这件装束也有其所能做的。\newline{}因为，首先从它穿戴的简便开始：它无需冗长的整理。因此，没有必要让任何工匠事先一天就正式地摆弄其褶皱的折痕，然后使其更精致优美，并将整块堆积的浮凸交给撑架看管；随后，在黎明时分，首先借助腰带束起那本应从一开始就织得适中长度的束腰外衣，再次审视浮凸，重新整理任何不平整，使一部分突出在左边，但（现在结束褶皱）从肩膀上向后拉出其环形部分，从而形成凹陷，并让右肩空闲，仍将其堆在左肩，为后背保留另一组类似的褶皱，就这样用一件负担来衣着人！\newline{}总之，我将不断地追问你自己的良心，你穿长袍时的第一感觉是什么？你觉得自己是被穿着，还是被背负？穿着衣服，还是扛着衣服？如果你回答否定，我将跟你回家；我将看到你跨过门槛后立即匆忙做什么。确实，没有一件衣服像长袍那样，脱下来让人感到庆幸。\newline{}我们不说鞋子——它们是长袍特有的折磨工具，对脚最不清洁的保护，是的，也是虚假的。因为在冷热中，谁不会觉得光脚僵硬比穿着鞋子束缚更便利呢？威尼斯的鞋厂以柔弱的靴子形式为脚步提供了强大的弹药！\newline{}但是，没有什么比斗篷更便捷的了，即使它像克拉特斯的斗篷那样是双层的。穿斗篷时绝不会强迫浪费时间，因为它的全部艺术在于宽松地遮盖。只需一次围裹，而且绝不失优雅，就能一下子完全遮盖人的每一部分。它要么露出肩膀，要么包住肩膀；在其他方面，它贴合肩膀；没有环绕的支撑，没有环绕的束带，不必担心褶皱保持位置是否可靠；它容易打理，容易重新调整；甚至脱下来也不需要挂在十字架上直到第二天。如果下面穿着衬衫，腰带的折磨就多余了；如果穿着鞋子一类的东西，那是最清洁的；否则，脚宁可光着——至少更男子气概，如果光着脚，比穿鞋子更男子气概。这些（我提出的理由）暂时是为斗篷的，就你以名字诋毁它而言。\newline{}然而，现在它以功能向你们挑战。\Quote{我，}它说，\Quote{对论坛、选举场或元老院不负任何责任；我不守谄媚的夜更，不占据讲坛，不在裁判官邸徘徊；我不染运河之气，不染格栅之气，不是长凳的常客，不是法律的肆意破坏者，不是吠叫的辩护者，不是法官，不是士兵，不是国王：我已远离民众。我唯一关心的是我自己：除了不关心之外，别无其他牵挂。你在隐居而非公开中更能享受更好的生活。但你会谴责我懒惰。确实，\InnerQuote{我们当为国家、帝国和产业而活。}这是古时之说。没有人是为他人而生的，因为注定为自己而死。无论如何，当我们来到\Person{伊壁鸠鲁}和\Person{芝诺}时，你将\InnerQuote{智者}的称号赋予整个\Emph{宁静}的导师群体，他们以\InnerQuote{至高}和\InnerQuote{独一}的快乐之名祝圣了那\Emph{宁静}。尽管如此，在某种程度上，甚至\Emph{我}也被允许为公众带来益处。从任何界石或祭坛，我习惯于开出道德的药物——这些药物在赋予公共事务、国家和帝国健康方面，将比\Emph{你们的}作品更为有效。事实上，如果我以赤刃与你们交锋，长袍对国家的伤害比胸甲更多。\newline{}此外，我不奉承任何恶习；我不宽恕任何麻木、任何懒惰的积垢。我用烧灼的铁对付那导致\Person{M. 图利乌斯}以超过4000英镑购买一张香橼木圆桌，以及\Person{阿西尼乌斯·加卢斯}以双倍价格购买一张同样的摩尔木普通桌（呵！他们多么看重木纹斑点！）的野心；或者再次，\Person{苏拉}制作重达一百磅的盘子。我担心那个天平太小，当一位\Person{德鲁西拉努斯}（而且他是\Person{克劳狄}的奴隶！）制作了一个重达500磅的托盘！——这个托盘对上述桌子或许是必不可少的，为它如果建了一个工场，就应该也建一个餐厅。同样，我将手术刀刺入导致\Person{维迪乌斯·波利奥}将奴隶投喂海鳗的不人道行为。的确，他高兴于他新奇的残暴，豢养了无牙、无爪、无角的陆地怪物：他乐意强迫他的鱼变成野兽，这些鱼（当然）要立即烹煮，以便在他自己的肚子里也能尝到他自己奴隶身体的滋味。我将割除导致演说家\Person{霍滕修斯}第一个有心为食物杀死孔雀的饕餮；导致\Person{奥菲迪乌斯·卢尔科}第一个用填料败坏肉类，并借助肉酱使它们产生不纯的味道；导致\Person{阿西尼乌斯·克勒尔}以近50英镑购买一条鲻鱼的肉；导致演员\Person{埃索普斯}在他的食品储藏室里保存一个价值近800英镑的盘子，由同样昂贵的鸟类组成，包括所有歌鸟和说话鸟；导致他的儿子，在尝过如此美味后，竟敢渴望更奢侈的东西：因为他吞下了珍珠——即便仅就其名而言也是昂贵的——我想是怕他吃得比他父亲更寒酸。我对\Person{尼禄}们、\Person{阿皮丘斯}们和\Person{鲁弗斯}们保持沉默。我将给\Person{斯考鲁斯}的不洁、\Person{库里乌斯}的赌博、\Person{安东尼}的放纵下一剂泻药。记住，这些我提到的人中，许多都是穿托加袍的人——这种人在穿斗篷的人中不易找到。这些国家的脓疮，除了披着斗篷的言论，谁能清除和排脓呢？}\par

% structure: p0026 source=h2 target=section reason=relative-heading-rank; base=h2

\section{第六章 披肩的进一步区分与最高荣耀}

\Quote{ \InnerQuote{凭言词，}（我的对手）说，\InnerQuote{你试图说服我——一个极其明智的药物。} 但是，尽管言语可能沉默——因幼年受阻或因羞怯而止，因为生命满足于甚至无舌的哲学——我的非常\Emph{裁剪}本身雄辩有力。事实上，一位哲学家被\Emph{听闻}，只要他被\Emph{看见}。我的目光使恶行羞愧。谁不痛苦，当他看到自己的对手？谁能忍受用眼睛凝视那在精神上他不能凝视的人？披肩所赐予的恩惠是伟大的，一想到它，道德败坏绝对羞愧。让哲学现在考虑她自己的利益问题；因为她不是我夸耀的唯一同伴。我夸耀其他有益于公众的科学技艺。来自我的储备，文字形式的第一位教师、声音的第一位解释者、算术基础的首位训练者、文法学家、修辞学家、智者、医生、诗人、音乐节拍者、占星学家和观鸟者都得到衣着。所有自由学艺都被我的四角所覆盖。\InnerQuote{确实；但所有这些地位都低于罗马骑士}好吧；但你的角斗士训练者及其所有可耻的随从，都穿着托加袍被带入竞技场。这无疑就是隐含在\InnerQuote{从长袍到披肩！}中的侮辱。} 好吧，\OriginalTerm{披肩}{Mantle}如此说。但我也授予它与一个神圣教派和训诫的共融。喜乐吧，披肩，欢庆吧！一种更好的哲学如今屈尊来尊敬你，因为你已开始成为基督徒的衣饰！\par

\SourceRecord{Translated by S. Thelwall.；Ante-Nicene Fathers；Vol. 4.；Edited by Alexander Roberts, James Donaldson, and A. Cleveland Coxe.；Buffalo, NY: Christian Literature Publishing Co.,；1885.；<http://www.newadvent.org/fathers/0401.htm>.}
